% Fichier généré par scripts/assemble_memoire.py : ne pas modifier à la main.
\usepackage{iftex}
\ifXeTeX\usepackage{fontspec}\else\usepackage[utf8]{inputenc}\usepackage[T1]{fontenc}\fi % XeTeX (tectonic) ou pdfLaTeX (Overleaf)
\usepackage[english,french]{babel}% français par défaut ; anglais pour l'abstract
\ifXeTeX\usepackage{microtype}\else\usepackage{lmodern,microtype}\fi
\usepackage[margin=2.5cm]{geometry}
\usepackage{booktabs,array,graphicx,xcolor,longtable}
\usepackage{tikz}
\usetikzlibrary{calc,arrows.meta,positioning}
\usepackage[round,authoryear]{natbib}
\usepackage[hidelinks]{hyperref}
\usepackage{setspace}\onehalfspacing
\setlength{\emergencystretch}{1.5em}% évite les lignes qui débordent (mots longs, citations)
\usepackage{eso-pic,amsmath,amssymb,placeins}
\usepackage{fvextra}% consignes des systèmes en annexe, avec retour à la ligne
\newcommand{\aremplir}[1]{\textbf{\textcolor{red}{[À remplir : #1]}}}

% --- Données fictives -------------------------------------------------------------------------
% genere/valeurs.tex est écrit par analysis/analyse_mini.py. Il bascule \fictiftrue si une seule ligne
% des fichiers JSONL porte "fictif": true. Avec les vrais runs, les marques disparaissent d'elles-mêmes.
\newif\iffictif
\newcommand{\fictifmarque}{}
% Généré par analysis/analyse_mini.py : ne pas modifier à la main.
\fictiftrue
\expandafter\def\csname valEM-A0-nq\endcsname{8}
\expandafter\def\csname valRech-A0-nq\endcsname{0{,}00}
\expandafter\def\csname valEpuise-A0-nq\endcsname{0}
\expandafter\def\csname valEM-A0-triviaqa\endcsname{39}
\expandafter\def\csname valRech-A0-triviaqa\endcsname{0{,}00}
\expandafter\def\csname valEpuise-A0-triviaqa\endcsname{0}
\expandafter\def\csname valEM-A1-nq\endcsname{18}
\expandafter\def\csname valRech-A1-nq\endcsname{1{,}00}
\expandafter\def\csname valEpuise-A1-nq\endcsname{0}
\expandafter\def\csname valEM-A1-triviaqa\endcsname{49}
\expandafter\def\csname valRech-A1-triviaqa\endcsname{1{,}00}
\expandafter\def\csname valEpuise-A1-triviaqa\endcsname{0}
\expandafter\def\csname valEM-A2-nq\endcsname{20}
\expandafter\def\csname valRech-A2-nq\endcsname{1{,}57}
\expandafter\def\csname valEpuise-A2-nq\endcsname{0}
\expandafter\def\csname valEM-A2-triviaqa\endcsname{41}
\expandafter\def\csname valRech-A2-triviaqa\endcsname{1{,}41}
\expandafter\def\csname valEpuise-A2-triviaqa\endcsname{0}
\expandafter\def\csname valEM-A3-nq\endcsname{20}
\expandafter\def\csname valRech-A3-nq\endcsname{3{,}27}
\expandafter\def\csname valEpuise-A3-nq\endcsname{4}
\expandafter\def\csname valEM-A3-triviaqa\endcsname{45}
\expandafter\def\csname valRech-A3-triviaqa\endcsname{3{,}33}
\expandafter\def\csname valEpuise-A3-triviaqa\endcsname{4}
\expandafter\def\csname valEM-A4-nq\endcsname{25}
\expandafter\def\csname valRech-A4-nq\endcsname{2{,}29}
\expandafter\def\csname valEpuise-A4-nq\endcsname{0}
\expandafter\def\csname valEM-A4-triviaqa\endcsname{43}
\expandafter\def\csname valRech-A4-triviaqa\endcsname{2{,}30}
\expandafter\def\csname valEpuise-A4-triviaqa\endcsname{2}
\expandafter\def\csname valEM-A5-nq\endcsname{24}
\expandafter\def\csname valRech-A5-nq\endcsname{1{,}95}
\expandafter\def\csname valEpuise-A5-nq\endcsname{1}
\expandafter\def\csname valEM-A5-triviaqa\endcsname{45}
\expandafter\def\csname valRech-A5-triviaqa\endcsname{1{,}94}
\expandafter\def\csname valEpuise-A5-triviaqa\endcsname{0}
\expandafter\def\csname valBest-rl-nq\endcsname{25}
\expandafter\def\csname valBestNom-rl-nq\endcsname{A4}
\expandafter\def\csname valBest-harnais-nq\endcsname{20}
\expandafter\def\csname valBestNom-harnais-nq\endcsname{A2}
\expandafter\def\csname valBest-rl-triviaqa\endcsname{45}
\expandafter\def\csname valBestNom-rl-triviaqa\endcsname{A5}
\expandafter\def\csname valBest-harnais-triviaqa\endcsname{49}
\expandafter\def\csname valBestNom-harnais-triviaqa\endcsname{A1}
\expandafter\def\csname valHA-HA1-nb_ok\endcsname{0}
\expandafter\def\csname valHA-HA1-verdict\endcsname{non soutenue}
\expandafter\def\csname valHA-HA2-d\endcsname{6{,}2}
\expandafter\def\csname valHA-HA2-lo\endcsname{$-$3{,}8}
\expandafter\def\csname valHA-HA2-hi\endcsname{15{,}7}
\expandafter\def\csname valHA-HA2-ecart_nq\endcsname{5{,}2}
\expandafter\def\csname valHA-HA2-ecart_tqa\endcsname{$-$1{,}0}
\expandafter\def\csname valHA-HA2-verdict\endcsname{non soutenue}
\expandafter\def\csname valHA-HA3a-d\endcsname{0{,}63}
\expandafter\def\csname valHA-HA3a-lo\endcsname{0{,}48}
\expandafter\def\csname valHA-HA3a-hi\endcsname{0{,}79}
\expandafter\def\csname valHA-HA3a-n\endcsname{200}
\expandafter\def\csname valHA-HA3a-verdict\endcsname{soutenue}
\expandafter\def\csname valHAen-HA3a-d\endcsname{0.63}
\expandafter\def\csname valHA-HA3b-d\endcsname{1{,}23}
\expandafter\def\csname valHA-HA3b-lo\endcsname{0{,}91}
\expandafter\def\csname valHA-HA3b-hi\endcsname{1{,}55}
\expandafter\def\csname valHA-HA3b-n\endcsname{47}
\expandafter\def\csname valHA-HA3b-verdict\endcsname{soutenue}
\expandafter\def\csname valHAen-HA3b-d\endcsname{1.23}

\newcommand{\valEM}[2]{\csname valEM-#1-#2\endcsname}         % EM en points, ex. \valEM{A4}{nq}
\newcommand{\valRech}[2]{\csname valRech-#1-#2\endcsname}     % recherches par question
\newcommand{\valHA}[2]{\csname valHA-#1-#2\endcsname}         % résultat d'un test, ex. \valHA{HA3a}{d}
\newcommand{\valHAen}[2]{\csname valHAen-#1-#2\endcsname}     % même valeur, point décimal (abstract anglais)
\newcommand{\valSansRech}[1]{\csname valSansRech-#1\endcsname} % part des questions sans recherche (%)
\newcommand{\valPerdues}[1]{\csname valPerdues-#1\endcsname}   % questions « connues » (A0 juste) ratées
\newcommand{\valSansRep}[1]{\csname valSansRep-#1\endcsname}   % sorties sans réponse (budget épuisé, format) sur 200
\newcommand{\valAbst}[1]{\csname valAbst-#1\endcsname}         % abstentions explicites sur 200
% genere/approf_valeurs.tex est écrit par analysis/analyse_approfondie.py (analyses exploratoires).
% Généré par analysis/analyse_approfondie.py : ne pas modifier à la main.
\expandafter\def\csname valAp-A1-trouvee\endcsname{37}
\expandafter\def\csname valAp-A1-em_t\endcsname{53}
\expandafter\def\csname valAp-A1-em_nt\endcsname{9}
\expandafter\def\csname valAp-A1-n_t\endcsname{74}
\expandafter\def\csname valAp-A1-n_nt\endcsname{126}
\expandafter\def\csname valAp-A1-trouvee1\endcsname{37}
\expandafter\def\csname valAp-A1-reform\endcsname{0}
\expandafter\def\csname valAp-A2-trouvee\endcsname{34}
\expandafter\def\csname valAp-A2-em_t\endcsname{46}
\expandafter\def\csname valAp-A2-em_nt\endcsname{8}
\expandafter\def\csname valAp-A2-n_t\endcsname{67}
\expandafter\def\csname valAp-A2-n_nt\endcsname{133}
\expandafter\def\csname valAp-A2-trouvee1\endcsname{31}
\expandafter\def\csname valAp-A2-reform\endcsname{2}
\expandafter\def\csname valAp-A3-trouvee\endcsname{42}
\expandafter\def\csname valAp-A3-em_t\endcsname{51}
\expandafter\def\csname valAp-A3-em_nt\endcsname{2}
\expandafter\def\csname valAp-A3-n_t\endcsname{85}
\expandafter\def\csname valAp-A3-n_nt\endcsname{115}
\expandafter\def\csname valAp-A3-trouvee1\endcsname{37}
\expandafter\def\csname valAp-A3-reform\endcsname{6}
\expandafter\def\csname valAp-A4-trouvee\endcsname{46}
\expandafter\def\csname valAp-A4-em_t\endcsname{67}
\expandafter\def\csname valAp-A4-em_nt\endcsname{7}
\expandafter\def\csname valAp-A4-n_t\endcsname{93}
\expandafter\def\csname valAp-A4-n_nt\endcsname{107}
\expandafter\def\csname valAp-A4-trouvee1\endcsname{36}
\expandafter\def\csname valAp-A4-reform\endcsname{11}
\expandafter\def\csname valAp-A5-trouvee\endcsname{38}
\expandafter\def\csname valAp-A5-em_t\endcsname{71}
\expandafter\def\csname valAp-A5-em_nt\endcsname{16}
\expandafter\def\csname valAp-A5-n_t\endcsname{77}
\expandafter\def\csname valAp-A5-n_nt\endcsname{123}
\expandafter\def\csname valAp-A5-trouvee1\endcsname{36}
\expandafter\def\csname valAp-A5-reform\endcsname{2}
\expandafter\def\csname valAp-harnais-trouvee\endcsname{38}
\expandafter\def\csname valAp-harnais-em_t\endcsname{50}
\expandafter\def\csname valAp-harnais-em_nt\endcsname{6}
\expandafter\def\csname valAp-harnais-n_t\endcsname{226}
\expandafter\def\csname valAp-harnais-n_nt\endcsname{374}
\expandafter\def\csname valAp-rl-trouvee\endcsname{42}
\expandafter\def\csname valAp-rl-em_t\endcsname{69}
\expandafter\def\csname valAp-rl-em_nt\endcsname{12}
\expandafter\def\csname valAp-rl-n_t\endcsname{170}
\expandafter\def\csname valAp-rl-n_nt\endcsname{230}
\expandafter\def\csname valAp-A1-gagnees\endcsname{24}
\expandafter\def\csname valAp-A1-perdues\endcsname{18}
\expandafter\def\csname valAp-A1-gardees\endcsname{26}
\expandafter\def\csname valAp-A2-gagnees\endcsname{20}
\expandafter\def\csname valAp-A2-perdues\endcsname{23}
\expandafter\def\csname valAp-A2-gardees\endcsname{21}
\expandafter\def\csname valAp-A3-gagnees\endcsname{25}
\expandafter\def\csname valAp-A3-perdues\endcsname{24}
\expandafter\def\csname valAp-A3-gardees\endcsname{20}
\expandafter\def\csname valAp-A4-gagnees\endcsname{32}
\expandafter\def\csname valAp-A4-perdues\endcsname{6}
\expandafter\def\csname valAp-A4-gardees\endcsname{38}
\expandafter\def\csname valAp-A5-gagnees\endcsname{41}
\expandafter\def\csname valAp-A5-perdues\endcsname{10}
\expandafter\def\csname valAp-A5-gardees\endcsname{34}
\expandafter\def\csname valAp-A1-copie1\endcsname{100}
\expandafter\def\csname valAp-A1-jacc\endcsname{1{,}00}
\expandafter\def\csname valAp-A1-repetees\endcsname{0}
\expandafter\def\csname valAp-A2-copie1\endcsname{4}
\expandafter\def\csname valAp-A2-jacc\endcsname{0{,}42}
\expandafter\def\csname valAp-A2-repetees\endcsname{6}
\expandafter\def\csname valAp-A3-copie1\endcsname{100}
\expandafter\def\csname valAp-A3-jacc\endcsname{0{,}45}
\expandafter\def\csname valAp-A3-repetees\endcsname{0}
\expandafter\def\csname valAp-A4-copie1\endcsname{70}
\expandafter\def\csname valAp-A4-jacc\endcsname{0{,}50}
\expandafter\def\csname valAp-A4-repetees\endcsname{29}
\expandafter\def\csname valAp-A5-copie1\endcsname{90}
\expandafter\def\csname valAp-A5-jacc\endcsname{0{,}85}
\expandafter\def\csname valAp-A5-repetees\endcsname{26}
\expandafter\def\csname valAp-A0-err0\endcsname{1}
\expandafter\def\csname valAp-A0-err1\endcsname{13}
\expandafter\def\csname valAp-A0-err2\endcsname{0}
\expandafter\def\csname valAp-A0-err3\endcsname{142}
\expandafter\def\csname valAp-A1-err0\endcsname{3}
\expandafter\def\csname valAp-A1-err1\endcsname{14}
\expandafter\def\csname valAp-A1-err2\endcsname{24}
\expandafter\def\csname valAp-A1-err3\endcsname{109}
\expandafter\def\csname valAp-A2-err0\endcsname{18}
\expandafter\def\csname valAp-A2-err1\endcsname{12}
\expandafter\def\csname valAp-A2-err2\endcsname{26}
\expandafter\def\csname valAp-A2-err3\endcsname{103}
\expandafter\def\csname valAp-A3-err0\endcsname{42}
\expandafter\def\csname valAp-A3-err1\endcsname{9}
\expandafter\def\csname valAp-A3-err2\endcsname{28}
\expandafter\def\csname valAp-A3-err3\endcsname{76}
\expandafter\def\csname valAp-A4-err0\endcsname{8}
\expandafter\def\csname valAp-A4-err1\endcsname{19}
\expandafter\def\csname valAp-A4-err2\endcsname{17}
\expandafter\def\csname valAp-A4-err3\endcsname{86}
\expandafter\def\csname valAp-A5-err0\endcsname{1}
\expandafter\def\csname valAp-A5-err1\endcsname{21}
\expandafter\def\csname valAp-A5-err2\endcsname{10}
\expandafter\def\csname valAp-A5-err3\endcsname{93}
\expandafter\def\csname valAp-A0-jetons\endcsname{11}
\expandafter\def\csname valAp-A1-jetons\endcsname{13}
\expandafter\def\csname valAp-A2-jetons\endcsname{151}
\expandafter\def\csname valAp-A3-jetons\endcsname{139}
\expandafter\def\csname valAp-A4-jetons\endcsname{162}
\expandafter\def\csname valAp-A5-jetons\endcsname{166}
\expandafter\def\csname valAp-ox-trouver\endcsname{2{,}1}
\expandafter\def\csname valAp-ox-trouver-lo\endcsname{0{,}9}
\expandafter\def\csname valAp-ox-trouver-hi\endcsname{3{,}4}
\expandafter\def\csname valAp-ox-trouver-part\endcsname{16}
\expandafter\def\csname valAp-ox-lire\endcsname{8{,}0}
\expandafter\def\csname valAp-ox-lire-lo\endcsname{4{,}9}
\expandafter\def\csname valAp-ox-lire-hi\endcsname{11{,}4}
\expandafter\def\csname valAp-ox-lire-part\endcsname{59}
\expandafter\def\csname valAp-ox-sans\endcsname{3{,}5}
\expandafter\def\csname valAp-ox-sans-lo\endcsname{0{,}7}
\expandafter\def\csname valAp-ox-sans-hi\endcsname{6{,}4}
\expandafter\def\csname valAp-ox-sans-part\endcsname{26}
\expandafter\def\csname valAp-ox-total\endcsname{13{,}6}
\expandafter\def\csname valAp-ox-total-lo\endcsname{9{,}2}
\expandafter\def\csname valAp-ox-total-hi\endcsname{18{,}2}

\newcommand{\valAp}[1]{\csname valAp-#1\endcsname}             % ex. \valAp{rl-em_t}, \valAp{ox-lire}
% genere/avance_valeurs.tex est écrit par analysis/analyse_avancee.py (analyses avancées, exploratoires).
% Généré par analysis/analyse_avancee.py : ne pas modifier à la main.
\expandafter\def\csname valAv-glm-rl\endcsname{2{,}62}
\expandafter\def\csname valAv-glm-rl-lo\endcsname{1{,}44}
\expandafter\def\csname valAv-glm-rl-hi\endcsname{4{,}78}
\expandafter\def\csname valAv-glm-rl-p\endcsname{0{,}002}
\expandafter\def\csname valAv-glm-trouvee\endcsname{12{,}96}
\expandafter\def\csname valAv-glm-trouvee-lo\endcsname{6{,}91}
\expandafter\def\csname valAv-glm-trouvee-hi\endcsname{24{,}31}
\expandafter\def\csname valAv-glm-trouvee-p\endcsname{$<$\,0{,}001}
\expandafter\def\csname valAv-glm-rlxtrouvee\endcsname{1{,}03}
\expandafter\def\csname valAv-glm-rlxtrouvee-lo\endcsname{0{,}50}
\expandafter\def\csname valAv-glm-rlxtrouvee-hi\endcsname{2{,}12}
\expandafter\def\csname valAv-glm-rlxtrouvee-p\endcsname{0{,}942}
\expandafter\def\csname valAv-glm-connue\endcsname{5{,}61}
\expandafter\def\csname valAv-glm-connue-lo\endcsname{2{,}87}
\expandafter\def\csname valAv-glm-connue-hi\endcsname{10{,}97}
\expandafter\def\csname valAv-glm-connue-p\endcsname{$<$\,0{,}001}
\expandafter\def\csname valAv-glm-tqa\endcsname{3{,}90}
\expandafter\def\csname valAv-glm-tqa-lo\endcsname{2{,}19}
\expandafter\def\csname valAv-glm-tqa-hi\endcsname{6{,}96}
\expandafter\def\csname valAv-glm-tqa-p\endcsname{$<$\,0{,}001}
\expandafter\def\csname valAv-glm-nobs\endcsname{1\,000}
\expandafter\def\csname valAv-glm-nq\endcsname{200}
\expandafter\def\csname valAv-A1-cum1\endcsname{37}
\expandafter\def\csname valAv-A1-marg1\endcsname{37{,}0}
\expandafter\def\csname valAv-A1-atteint1\endcsname{100}
\expandafter\def\csname valAv-A1-cum2\endcsname{37}
\expandafter\def\csname valAv-A1-marg2\endcsname{0{,}0}
\expandafter\def\csname valAv-A1-atteint2\endcsname{0}
\expandafter\def\csname valAv-A1-cum3\endcsname{37}
\expandafter\def\csname valAv-A1-marg3\endcsname{0{,}0}
\expandafter\def\csname valAv-A1-atteint3\endcsname{0}
\expandafter\def\csname valAv-A1-cum4\endcsname{37}
\expandafter\def\csname valAv-A1-marg4\endcsname{0{,}0}
\expandafter\def\csname valAv-A1-atteint4\endcsname{0}
\expandafter\def\csname valAv-A2-cum1\endcsname{31}
\expandafter\def\csname valAv-A2-marg1\endcsname{31{,}0}
\expandafter\def\csname valAv-A2-atteint1\endcsname{96}
\expandafter\def\csname valAv-A2-cum2\endcsname{34}
\expandafter\def\csname valAv-A2-marg2\endcsname{2{,}5}
\expandafter\def\csname valAv-A2-atteint2\endcsname{29}
\expandafter\def\csname valAv-A2-cum3\endcsname{34}
\expandafter\def\csname valAv-A2-marg3\endcsname{0{,}0}
\expandafter\def\csname valAv-A2-atteint3\endcsname{10}
\expandafter\def\csname valAv-A2-cum4\endcsname{34}
\expandafter\def\csname valAv-A2-marg4\endcsname{0{,}0}
\expandafter\def\csname valAv-A2-atteint4\endcsname{6}
\expandafter\def\csname valAv-A3-cum1\endcsname{37}
\expandafter\def\csname valAv-A3-marg1\endcsname{37{,}0}
\expandafter\def\csname valAv-A3-atteint1\endcsname{100}
\expandafter\def\csname valAv-A3-cum2\endcsname{40}
\expandafter\def\csname valAv-A3-marg2\endcsname{3{,}0}
\expandafter\def\csname valAv-A3-atteint2\endcsname{100}
\expandafter\def\csname valAv-A3-cum3\endcsname{41}
\expandafter\def\csname valAv-A3-marg3\endcsname{1{,}0}
\expandafter\def\csname valAv-A3-atteint3\endcsname{80}
\expandafter\def\csname valAv-A3-cum4\endcsname{42}
\expandafter\def\csname valAv-A3-marg4\endcsname{1{,}5}
\expandafter\def\csname valAv-A3-atteint4\endcsname{35}
\expandafter\def\csname valAv-A4-cum1\endcsname{36}
\expandafter\def\csname valAv-A4-marg1\endcsname{35{,}5}
\expandafter\def\csname valAv-A4-atteint1\endcsname{100}
\expandafter\def\csname valAv-A4-cum2\endcsname{46}
\expandafter\def\csname valAv-A4-marg2\endcsname{10{,}0}
\expandafter\def\csname valAv-A4-atteint2\endcsname{100}
\expandafter\def\csname valAv-A4-cum3\endcsname{46}
\expandafter\def\csname valAv-A4-marg3\endcsname{1{,}0}
\expandafter\def\csname valAv-A4-atteint3\endcsname{66}
\expandafter\def\csname valAv-A4-cum4\endcsname{46}
\expandafter\def\csname valAv-A4-marg4\endcsname{0{,}0}
\expandafter\def\csname valAv-A4-atteint4\endcsname{22}
\expandafter\def\csname valAv-A5-cum1\endcsname{36}
\expandafter\def\csname valAv-A5-marg1\endcsname{36{,}5}
\expandafter\def\csname valAv-A5-atteint1\endcsname{100}
\expandafter\def\csname valAv-A5-cum2\endcsname{38}
\expandafter\def\csname valAv-A5-marg2\endcsname{1{,}0}
\expandafter\def\csname valAv-A5-atteint2\endcsname{68}
\expandafter\def\csname valAv-A5-cum3\endcsname{38}
\expandafter\def\csname valAv-A5-marg3\endcsname{1{,}0}
\expandafter\def\csname valAv-A5-atteint3\endcsname{10}
\expandafter\def\csname valAv-A5-cum4\endcsname{38}
\expandafter\def\csname valAv-A5-marg4\endcsname{0{,}0}
\expandafter\def\csname valAv-A5-atteint4\endcsname{0}
\expandafter\def\csname valAv-sens-em\endcsname{13{,}6}
\expandafter\def\csname valAv-sens-em-lo\endcsname{9{,}2}
\expandafter\def\csname valAv-sens-em-hi\endcsname{18{,}2}
\expandafter\def\csname valAv-sens-em-nb\endcsname{6}
\expandafter\def\csname valAv-sens-souple\endcsname{12{,}5}
\expandafter\def\csname valAv-sens-souple-lo\endcsname{8{,}1}
\expandafter\def\csname valAv-sens-souple-hi\endcsname{17{,}2}
\expandafter\def\csname valAv-sens-souple-nb\endcsname{6}
\expandafter\def\csname valAv-sens-f1\endcsname{15{,}5}
\expandafter\def\csname valAv-sens-f1-lo\endcsname{11{,}3}
\expandafter\def\csname valAv-sens-f1-hi\endcsname{19{,}9}
\expandafter\def\csname valAv-or-min\endcsname{3{,}8}
\expandafter\def\csname valAv-or-max\endcsname{6{,}0}
\expandafter\def\csname valAv-jeu-nq-trouver\endcsname{1{,}2}
\expandafter\def\csname valAv-jeu-nq-lire\endcsname{6{,}5}
\expandafter\def\csname valAv-jeu-nq-sans\endcsname{5{,}1}
\expandafter\def\csname valAv-jeu-nq-total\endcsname{12{,}8}
\expandafter\def\csname valAv-jeu-triviaqa-trouver\endcsname{2{,}7}
\expandafter\def\csname valAv-jeu-triviaqa-lire\endcsname{9{,}7}
\expandafter\def\csname valAv-jeu-triviaqa-sans\endcsname{2{,}0}
\expandafter\def\csname valAv-jeu-triviaqa-total\endcsname{14{,}3}
\expandafter\def\csname valAv-comp-oracle_h\endcsname{34}
\expandafter\def\csname valAv-comp-oracle_rl\endcsname{43}
\expandafter\def\csname valAv-comp-oracle_tous\endcsname{48}
\expandafter\def\csname valAv-comp-jac_rl\endcsname{0{,}69}
\expandafter\def\csname valAv-comp-jac_h\endcsname{0{,}48}
\expandafter\def\csname valAv-comp-aucun\endcsname{52}

\newcommand{\valAv}[1]{\csname valAv-#1\endcsname}             % ex. \valAv{glm-rl}, \valAv{A4-marg2}
\newcommand{\valBest}[2]{\csname valBest-#1-#2\endcsname}     % EM du meilleur système d'une famille (rl, harnais)
\newcommand{\valBestNom}[2]{\csname valBestNom-#1-#2\endcsname}
% genere/puissance_valeurs.tex est écrit par analysis/puissance.py (simulation, indépendante des données).
% Généré par analysis/puissance.py
\expandafter\def\csname valPuiss-100-5-simple\endcsname{13\,\%}
\expandafter\def\csname valPuiss-100-5-holm\endcsname{3\,\%}
\expandafter\def\csname valPuiss-100-10-simple\endcsname{42\,\%}
\expandafter\def\csname valPuiss-100-10-holm\endcsname{19\,\%}
\expandafter\def\csname valPuiss-100-15-simple\endcsname{81\,\%}
\expandafter\def\csname valPuiss-100-15-holm\endcsname{56\,\%}
\expandafter\def\csname valPuiss-200-5-simple\endcsname{28\,\%}
\expandafter\def\csname valPuiss-200-5-holm\endcsname{10\,\%}
\expandafter\def\csname valPuiss-200-10-simple\endcsname{79\,\%}
\expandafter\def\csname valPuiss-200-10-holm\endcsname{54\,\%}
\expandafter\def\csname valPuiss-200-15-simple\endcsname{99\,\%}
\expandafter\def\csname valPuiss-200-15-holm\endcsname{95\,\%}
\expandafter\def\csname valPuiss-500-5-simple\endcsname{65\,\%}
\expandafter\def\csname valPuiss-500-5-holm\endcsname{35\,\%}
\expandafter\def\csname valPuiss-500-10-simple\endcsname{99\,\%}
\expandafter\def\csname valPuiss-500-10-holm\endcsname{96\,\%}
\expandafter\def\csname valPuiss-500-15-simple\endcsname{100\,\%}
\expandafter\def\csname valPuiss-500-15-holm\endcsname{100\,\%}
\expandafter\def\csname valPuiss-1000-5-simple\endcsname{90\,\%}
\expandafter\def\csname valPuiss-1000-5-holm\endcsname{74\,\%}
\expandafter\def\csname valPuiss-1000-10-simple\endcsname{100\,\%}
\expandafter\def\csname valPuiss-1000-10-holm\endcsname{100\,\%}
\expandafter\def\csname valPuiss-1000-15-simple\endcsname{100\,\%}
\expandafter\def\csname valPuiss-1000-15-holm\endcsname{100\,\%}

\newcommand{\valPuiss}[3]{\csname valPuiss-#1-#2-#3\endcsname} % ex. \valPuiss{200}{10}{holm}
\newcommand{\fictif}{\iffictif\textcolor{red}{\textsuperscript{\,fictif}}\fi}
\newcommand{\fictifnote}[1]{\iffictif\par\noindent\fcolorbox{red}{red!4}{\parbox{\dimexpr\linewidth-2\fboxsep-2\fboxrule}{\small\textcolor{red!70!black}{\textbf{Données fictives.} #1}}}\par\medskip\fi}
\iffictif
\AddToShipoutPictureBG{\AtPageCenter{\makebox(0,0){\rotatebox{55}{\textcolor{red!12}{\fontsize{60}{60}\selectfont\bfseries DONNÉES FICTIVES}}}}}
\fi

